\section{Causal simulation and a joint attained-error law}
\begin{proposition}[Complete common-task composition]\label{prop:simulation}
A simulator from a source experiment to a target experiment must specify a parameter-independent strategy lift, report transformation, internal state, clock, task readout and physical-command compatibility. Suppose it uses at most $S$ retained states, a target observer uses $M$, and their entire stopped scored transcript laws differ by at most $\eps$ in total variation, with convention $\|P-Q\|_{\TV}=\sup_A|P(A)-Q(A)|$. For loss in $[0,L]$, the composed observer uses at most $SM$ states and its risk differs by at most $L\eps$. Its expected cost bound requires its own additive-cost or killed-potential certificate. Two compatible causal simulators have total deficiency at most the sum of their deficiencies and a product state upper bound.
\end{proposition}
\begin{proof}
The simulator and observer retain their ordered pair of states; clocks and any continuing random seeds are part of those states. Compatibility makes the composition a legal causal procedure. For a common bounded scored loss, the layer-cake formula and the total-variation convention give the risk bound. Contractivity of total variation under a common Markov kernel and the triangle inequality give the two-simulator bound on complete stopped scored laws. Nonnegative costs must instead be summed along actual executed calls or bounded by separately certified stopped potentials. A small terminal-law discrepancy alone can coexist with arbitrarily large expected waiting. Reverse risk bounds need a reverse simulator or a separate attained lower construction.
\end{proof}

The state product is an implementation upper, not a claim of optimal simulator size. Root-mean-square readout errors add by the triangle inequality before squaring; deterministic errors may alternatively use the bounded-loss estimate following Proposition~\ref{prop:relative}. Neither convention makes an allowed upper bound on calibration or simulation error an unavoidable risk lower.


\begin{proposition}[Composition with live--terminal overlap]\label{prop:role-composition}
Suppose a target program has $\ell$ live roles and $t$ terminal roles. A compatible parameter-independent simulator has at most $K$ internal states while the target is live, including every continuing phase, seed and clock. Suppose further that upon target commitment all information needed for its final exported readout is represented by one of $E$ terminal export classes; all other live simulator state may then be discarded. The composed program has a realization with at most $\ell K$ live and $tE$ terminal roles. It therefore needs at most
\[
 \ell K+tE\quad\hbox{state-stopped labels},\qquad
 \max\{\ell K,tE\}\quad\hbox{event-stopped labels}.
\]
Under a complete stopped/scored-law deficiency certificate $\eps$, its bounded-loss discrepancy is at most $L\eps$. For two such simulators, $(K_1,E_1)$ and $(K_2,E_2)$ combine to $(K_1K_2,E_1E_2)$, with additive complete-law deficiencies. Without a terminal sufficient-class certificate, retain the full simulator state and use the safe product construction of Proposition~\ref{prop:simulation}.
\end{proposition}
\begin{proof}
Before commitment store the pair of the target live role and simulator internal state. Compatibility supplies its legal command distribution and post-report transitions, with no free port input. On commitment encode only the target output role and the sufficient export class, whose fixed decoder is part of the program. These are $\ell K$ live and $tE$ terminal roles. The complete stopped transcript, including the commitment event, is preserved by Lemma~\ref{lem:packing}; apply it to those two sets. A second composition stores the corresponding triples, giving the two products. Associativity is ordinary rebracketing of finite pairs and does not assert minimality. The complete-law risk and deficiency assertions follow exactly as in Proposition~\ref{prop:simulation}. Expected time and physical-call cost require their separate certificates at each composition.
\end{proof}

The compared transcript includes the mapped stopping time, the event of premature commitment, all required report exports and the audit. Stopping the source before the target is allowed only to the extent covered by that complete-law error certificate; a terminal-output marginal bound cannot justify it. A stateful translator is not a free change of notation: its continuing and terminal information appear in $K,E$. Thus this is a resource-certified simulation law, not an assertion of invariant scalar minimal memory under arbitrary simulators.

\begin{theorem}[Same-task attained calibration and erasure]\label{thm:joint}
Use the exact balanced class $D_0=1$, $B=2$, with matched effort. Add an initial paid observation of an independent fair sign $E_0$ erased with probability $2\eps$, where $0\le\eps\le1/2$, and a final paid audit of
\[
 (Z,E_0,s\delta),\qquad s\in\{-1,1\},\quad\delta\ge0.
\]
The fixed unknown sign $s$ affects no pre-audit kernel. Conditional on the initial reported tag, the entire core experiment contains no additional information about $E_0$ on the erased branch; its physical kernels may not read that hidden sign. All output coordinates are committed before this audit. Use squared Euclidean loss on the product carrier. For $N\ge\kappa+2$, put $b=(N-2)/\kappa$. The following bounds hold for either stopping interface $\diamond$ on its indicated sufficient state range:
\begin{align}
 \Qm_M(b)+2\eps+\delta^2
 &\le R_{{\rm joint},\diamond}^\Theta(N,M)
 \le\Qm_{m_\diamond}(b)+2\eps+\delta^2,\label{eq:joint-bounds}\\
 m_\st&=\lfloor(M-1)/3\rfloor-2\quad(M\ge13),\qquad
 m_\evt=\lfloor M/3\rfloor\quad(M\ge7).\label{eq:joint-widths}
\end{align}
If \eqref{eq:rectangle} holds, this gives in both cases
\begin{equation}\label{eq:joint-rate}
 R_{{\rm joint},\diamond}^\Theta(N,M)
 \asymp F((N-2)/\kappa,M)+\delta^2+\eps.
\end{equation}
The early report channel has causal deficiency \emph{exactly} $\eps$ relative to the perfect-sign channel when its virtual sign must be delivered before any core call. The stated alphabet ranges are sufficient constructions, not minimal-state assertions for the enlarged task.
\end{theorem}
\begin{proof}
For the upper use one initial live role and three tagged copies of the two core live roles: tags $+$, $-$ and erasure. Thus there are seven live roles. There are $3m$ terminal roles, one for each tag and each core output. Output the reported sign on known branches, zero on erasure, and zero for the inaccessible calibration sign. The balanced core upper has the same risk after every preceding history, so its contribution is $\Qm_m(b)$. The other contributions are exactly $2\eps$ and $\delta^2$. Initial observation and final audit each cost one. State stopping uses $7+3m=1+3(2+m)$ labels, whereas event stopping uses $\max\{7,3m\}$ by Lemma~\ref{lem:packing}. These give \eqref{eq:joint-widths}. The initial live role is distinct in this construction; no argument claims that every implementation needs such a separate label.

For the lower use the fixed reference core. Expected core cost is at most $N-2$, and its terminal projected dictionary has at most $M$ values under either semantics. Lemma~\ref{lem:lower} gives $\Qm_M(b)$. On the erased branch, induction over the causal core kernels and independent observer coins leaves $E_0$ fair conditional on all pre-audit information, even after stopping. Any sign estimate then has conditional squared loss at least one. Average the two indistinguishable alternatives $s=\pm1$: for any committed last coordinate $r$, their average loss is $\delta^2+r^2\ge\delta^2$. The pre-audit law, the core lower and the erased-sign lower are identical for these alternatives. The three lower contributions therefore hold simultaneously in a single worst-case total risk, not for incompatible separately optimized tasks. On the stated ranges, $m_\st\ge M/8$ and $m_\evt\ge M/4$. Apply Proposition~\ref{prop:geometry} to obtain \eqref{eq:joint-rate}.

On erasure, every immediately supplied virtual sign is independent of the true fair sign. The perfect channel has event \{virtual sign equals $E_0$\} with probability one, while every source simulator has probability at most $1-\eps$. This gives deficiency at least $\eps$ on the complete scored transcript. Reporting the known sign when available and an independent fair sign otherwise disagrees with the perfect report with probability $\eps$. A coupling, followed by the same causal core strategy and task audit, gives the opposite bound. The early completion requirement forbids waiting for the final audit to learn $E_0$.
\end{proof}

The joint theorem is an attained construction for one specified scored task, not a universal additive theorem for arbitrary error tolerances. Theorem~\ref{thm:approx} and Proposition~\ref{prop:relative} add their own budget correction and numerical certificates. Persistent labels, simulator states, terminal export classes, scratch, program description, pulse precision, raw acquisitions and physical time remain distinct accounts.
